% !TEX root = ../../main.tex
Intisari ditulis menggunakan bahasa Indonesia dengan jarak antar baris 1 spasi dan maksimal 1 halaman. Intisari sekurang-kurangnya berisi tentang latar belakang dan tujuan penelitian, metodologi yang digunakan, hasil penelitian, kesimpulan dan implikasi, dan Kata kunci yang berhubungan dengan penelitian.

Kata Kunci ditulis maksimal 5 kata yang paling berhubungan dengan isi skripsi. Silakan mengacu pada ACM / IEEE \textit{Computing classification} jika Anda adalah mahasiswa Sarjana TI \textcolor{blue}{http://www.acm.org/about/class/} atau mengacu kepada IEEE keywords \textcolor{blue}{http://www.ieee.org/documents/taxonomy\_v101.pdf} jika Anda berasal dari Prodi Sarjana TE.

\noindent{Kata kunci} : Kata kunci 1, Kata kunci 2, Kata kunci 3, Kata kunci 4, Kata kunci 5

\vspace{1cm}

%HAPUS YANG TIDAK PERLU
%-------------------------------------------------
\noindent\fbox{%
	\parbox{\textwidth}{%
\textbf{Contoh Abstrak Teknik Elektro:} \\

\hspace{1cm} "Penelitian ini bertujuan untuk mengembangkan sistem pengendalian suhu ruangan dengan menggunakan microcontroller. Metodologi yang digunakan adalah desain sirkuit, implementasi sistem pengendalian, dan pengujian performa. Hasil penelitian menunjukkan 
bahwa sistem pengendalian suhu ruangan yang dikembangkan mampu mengendalikan suhu ruangan dengan akurasi sebesar ±0,5°C. Kesimpulan dari penelitian ini adalah sistem pengendalian suhu ruangan yang dikembangkan efektif dan efisien. \\

Kata kunci: microcontroller, sistem pengendalian suhu, akurasi."
\vspace{5mm}


	}%
}

%-------------------------------------------------
